\section{Smoothed prime sums and nonvanishing}
\label{sec:prime-sum-reduction}

The purpose of this section is to identify a precise analytic estimate
which suffices for nonvanishing.  A global upper bound for a certain real
combination of smoothed prime sums is sufficient; a corresponding lower
bound is not required.  This is a conditional reduction.  No such upper
bound is established here.

\subsection{Normalization and Mellin transforms}

Let $\Lambda(n)$ denote the von Mangoldt function: it is $\log p$ when
$n=p^k$ for a prime $p$ and an integer $k\geq1$, and is zero otherwise.
Let $\gamma=\lim_{M\to\infty}(\sum_{n=1}^M n^{-1}-\log M)$.
For $x\geq1$, define
\begin{align}
 S(x)&=\sum_{n\leq x}\Lambda(n)(x/n-1),\label{mellin:S}\\
 G(x)&=S(x)-x\log x+(1+\gamma)x.\label{mellin:G}
\end{align}
For a nonprincipal Dirichlet character $\chi$, assume explicitly that
$L(1,\chi)\neq0$, and set
\begin{align}
 c_\chi&=\frac{L'(1,\chi)}{L(1,\chi)},\label{mellin:cchi}\\
 g_\chi(x)&=\sum_{n\leq x}\Lambda(n)\chi(n)(x/n-1)+c_\chi x,
 \qquad G_\chi(x)=\operatorname{Re}g_\chi(x).\label{mellin:gchi}
\end{align}
The assumption at $1$ is the classical nonvanishing theorem for
nonprincipal Dirichlet characters.  In the conditional statements below
it is an explicit hypothesis.  Each function in
\eqref{mellin:S}--\eqref{mellin:gchi} is continuous on $[1,\infty)$:
the newly appearing term at an integer $n$ has value zero at that point.
The elementary inequalities $|\chi(n)|\leq1$ and
$0\leq\Lambda(n)\leq\log n$ give
\begin{equation}\label{mellin:trivial-growth}
 |G(x)|+|g_\chi(x)|=O_\chi\bigl(x\log^2(2x)\bigr).
\end{equation}
For example, the uncorrected sum is bounded in absolute value by
$x\sum_{n\leq x}(\log n)/n$.

We record explicitly the analytic information needed near $s=1$.
For $\operatorname{Re}s>1$, partial summation gives
\begin{equation}\label{mellin:zeta-fractional}
 \zeta(s)=\frac{s}{s-1}
       -s\int_1^\infty\{x\}x^{-s-1}\,dx,
\end{equation}
where $\{x\}=x-\lfloor x\rfloor$.  Since $0\leq\{x\}<1$, the
integral is holomorphic for $\operatorname{Re}s>0$, by locally uniform
absolute convergence, also after differentiation.  This proves a
meromorphic continuation of $\zeta$ to that half-plane with just a simple
pole at $1$.  At an integer $M\geq1$,
\[
 \int_1^M\frac{\{x\}}{x^2}\,dx
   =\log M-\sum_{n=1}^{M-1}n\left(\frac1n-\frac1{n+1}\right)
   =1+\log M-\sum_{n=1}^M\frac1n.
\]
Consequently the constant term of~\eqref{mellin:zeta-fractional} at
$1$ is $\gamma$, and
\begin{equation}\label{mellin:zeta-laurent}
 \zeta(s)=\frac1{s-1}+\gamma+O(s-1),\qquad
 Z(s):=(s-1)\zeta(s),\quad Z(1)=1,\quad Z'(1)=\gamma.
\end{equation}
The harmonic limit defining $\gamma$ exists by integral comparison;
$\sum_{n=1}^M n^{-1}-\log M$ is positive and decreasing.

For a nonprincipal character, the sum of its values in a complete period
is zero.  Indeed, multiplication by a unit on which the character is
not $1$ permutes that sum and multiplies it by a different scalar.
Thus $A_\chi(x)=\sum_{n\leq x}\chi(n)$ is bounded.  The formula
\begin{equation}\label{mellin:dirichlet-continuation}
 L(s,\chi)=s\int_1^\infty A_\chi(x)x^{-s-1}\,dx
\end{equation}
continues $L(s,\chi)$ holomorphically to $\operatorname{Re}s>0$.
The equality first follows by summation by parts for real part greater
than $1$, and the integral converges locally uniformly for real part
greater than $0$.  These elementary continuations agree with the usual
$L$-functions by uniqueness of analytic continuation.

\begin{lemma}\label{mellin:identities}
For $\operatorname{Re}s>1$ the following integrals converge absolutely:
\begin{align}
 \int_1^\infty G(x)x^{-s-1}\,dx
 &=\frac{-\zeta'(s)/\zeta(s)}{s(s-1)}
     -\frac1{(s-1)^2}+\frac{1+\gamma}{s-1},\label{mellin:G-transform}\\
 \int_1^\infty g_\chi(x)x^{-s-1}\,dx
 &=\frac{-L'(s,\chi)/L(s,\chi)}{s(s-1)}
     +\frac{c_\chi}{s-1}.\label{mellin:g-transform}
\end{align}
The expression in~\eqref{mellin:G-transform} has a removable singularity
at $1$.  The same is true of~\eqref{mellin:g-transform} when
$L(1,\chi)\neq0$.
\end{lemma}

\begin{proof}
For $n\geq1$ and $\operatorname{Re}s>1$, direct integration gives
\[
 \int_n^\infty(x/n-1)x^{-s-1}\,dx
   =\frac{n^{-s}}{s(s-1)}.
\]
The sum of the absolute values of the corresponding integrals, with
coefficients bounded by $\Lambda(n)$, is finite.  This follows from
$\sum_{n\geq2}(\log n)n^{-\operatorname{Re}s}<\infty$.
Fubini's theorem therefore permits the exchange of the sum and the
integral.  The absolutely convergent Euler products give
\[
 -\frac{\zeta'(s)}{\zeta(s)}
   =\sum_{n\geq1}\frac{\Lambda(n)}{n^s},\qquad
 -\frac{L'(s,\chi)}{L(s,\chi)}
   =\sum_{n\geq1}\frac{\Lambda(n)\chi(n)}{n^s}.
\]
One obtains these identities by expanding the logarithm of each Euler
factor and differentiating its locally uniformly convergent series.
The same absolute convergence proves that the Euler products are
nonzero in this half-plane.
Finally,
\[
 \int_1^\infty x\,x^{-s-1}\,dx=\frac1{s-1},\qquad
 \int_1^\infty x\log x\,x^{-s-1}\,dx=\frac1{(s-1)^2}.
\]
These calculations prove both identities.

For the assertion at $1$, write $u=s-1$.  Equation
\eqref{mellin:zeta-laurent} implies
$-\zeta'/\zeta=1/u-\gamma+O(u)$.  Hence
\[
 \frac{-\zeta'/\zeta}{s(s-1)}
   =\frac1{u^2}-\frac{1+\gamma}{u}+O(1),
\]
so the two displayed correction terms cancel both polar terms.
For $\chi$, the quotient $L'/L$ is analytic near $1$ with value
$c_\chi$.  The polar term in its negative divided by $s(s-1)$ is
$-c_\chi/(s-1)$, again exactly cancelled by the correction term.
\end{proof}

The coefficient $1+\gamma$ and the sign of $c_\chi$ in these definitions
are therefore forced by cancellation at $1$.  Altering either linear
term changes the problem being estimated.

\subsection{The real singularity of a nonnegative Mellin transform}

The following version of Landau's principle supplies the analytic
reason that a one-sided bound can be enough.

\begin{lemma}\label{mellin:landau}
Let $H:[1,\infty)\to[0,\infty)$ be measurable and locally integrable,
and suppose $\int_1^\infty H(x)x^{-\sigma-1}\,dx$ is finite for at least
one real $\sigma$.  Put
\[
 \sigma_c=\inf\left\{\sigma\in\mathbb R:
       \int_1^\infty H(x)x^{-\sigma-1}\,dx<\infty\right\}.
\]
If $\sigma_c$ is finite, the holomorphic function
\[
 F(s)=\int_1^\infty H(x)x^{-s-1}\,dx,
       \qquad \operatorname{Re}s>\sigma_c,
\]
has no holomorphic continuation to a neighborhood of the real point
$\sigma_c$.
\end{lemma}

\begin{proof}
Nonnegativity makes the set of real exponents for which the integral
is finite an upper interval.  For every $\sigma>\sigma_c$, one can
choose a smaller exponent in that interval; hence the integral is
finite at $\sigma$.  Local uniform domination then makes $F$ holomorphic
in the indicated half-plane.  For a real $a>\sigma_c$,
\begin{equation}\label{mellin:moments}
 F^{(k)}(a)=(-1)^k\int_1^\infty
          H(x)(\log x)^k x^{-a-1}\,dx.
\end{equation}
Differentiation is justified by comparison with a slightly smaller
exponent greater than $\sigma_c$: every fixed power of $\log x$ is
bounded by a constant times an arbitrarily small positive power of $x$.

Suppose now that $F$ continues holomorphically to a disk
$|s-\sigma_c|<r$, for some $r>0$.  Set
$a=\sigma_c+r/4$ and $h=r/2$.  The Taylor series of this continuation
about $a$ converges at $a-h$, since the disk about $a$ of radius $3r/4$
is contained in the first disk.  Using~\eqref{mellin:moments}, all the
terms in this Taylor expansion are nonnegative real numbers:
\begin{align*}
 F(a-h)
  &=\sum_{k=0}^{\infty}\frac{(-h)^kF^{(k)}(a)}{k!}\\
  &=\sum_{k=0}^{\infty}\frac{h^k}{k!}
       \int_1^\infty H(x)(\log x)^k x^{-a-1}\,dx\\
  &=\int_1^\infty H(x)x^{-a-1}
       \sum_{k=0}^{\infty}\frac{(h\log x)^k}{k!}\,dx\\
  &=\int_1^\infty H(x)x^{-(a-h)-1}\,dx.
\end{align*}
The exchange uses Tonelli's theorem, so no integrability assumption
below $\sigma_c$ has been inserted.  The finite Taylor sum limit proves
that the last integral is finite.  But
$a-h=\sigma_c-r/4<\sigma_c$, contradicting the definition of
$\sigma_c$.
\end{proof}

\subsection{A global one-sided criterion for a family}

Let $\mathcal X$ be a finite set of nonprincipal Dirichlet characters,
possibly of different moduli.  Suppose that $\mathcal X$ is closed
under complex conjugation, and assign positive integer multiplicities
$m_\chi$ satisfying $m_{\overline\chi}=m_\chi$.  Assume
$L(1,\chi)\neq0$ for every $\chi\in\mathcal X$.  Let $R$ be a positive
integer and define the real-valued function
\begin{equation}\label{mellin:B}
 B(x)=R G(x)+\sum_{\chi\in\mathcal X}m_\chi G_\chi(x).
\end{equation}
The empty family is allowed.  It yields a criterion for $\zeta$ alone.

\begin{theorem}\label{mellin:one-sided-family}
Let $0\leq\theta<1$.  If there are constants $C_0\in\mathbb R$ and
$x_0\geq1$ such that
\begin{equation}\label{mellin:global-bound}
 B(x)\leq C_0x^\theta\qquad\text{for every }x\geq x_0,
\end{equation}
then $\zeta$ has no zeros in $\operatorname{Re}s>\theta$, and
$L(s,\chi)\neq0$ in this half-plane for every
$\chi\in\mathcal X$.  The pole of $\zeta$ at $1$ is understood as a
pole, not as a value to which the assertion is applied.
\end{theorem}

\begin{proof}
The identity
$L(s,\overline\chi)=\overline{L(\overline s,\chi)}$ follows first
from the Dirichlet series and then by analytic continuation.  In
particular $c_{\overline\chi}=\overline{c_\chi}$.  Therefore
\[
 \sum_{\chi\in\mathcal X}m_\chi G_\chi(x)
   =\sum_{\chi\in\mathcal X}m_\chi g_\chi(x),
 \qquad
 c=\sum_{\chi\in\mathcal X}m_\chi c_\chi\in\mathbb R.
\]
In the half-plane $\operatorname{Re}s>0$, form the holomorphic function
\[
 P(s)=Z(s)^R\prod_{\chi\in\mathcal X}L(s,\chi)^{m_\chi}.
\]
It is not identically zero, since the Euler products are nonzero for
$\operatorname{Re}s>1$.  Also $P(1)\neq0$ by the assumptions at $1$.
Put
\[
 a=-R,\qquad b=R(1+\gamma)+c.
\]
Logarithmic differentiation at $1$ gives
\begin{equation}\label{mellin:P-at-one}
 \frac{P'(1)}{P(1)}=R\gamma+c=a+b.
\end{equation}
Adding the identities in Lemma~\ref{mellin:identities} and using
$Z'/Z=1/(s-1)+\zeta'/\zeta$ yields, for
$\operatorname{Re}s>1$,
\begin{equation}\label{mellin:B-transform}
 \int_1^\infty B(x)x^{-s-1}\,dx
   =\frac{-P'(s)/P(s)+a+bs}{s(s-1)}.
\end{equation}

By continuity on $[1,x_0]$, increase $C_0$ to a nonnegative constant
$C$ so that
$H(x)=Cx^\theta-B(x)\geq0$ for all $x\geq1$.
Equation~\eqref{mellin:trivial-growth} shows that the Mellin transform
of $H$ converges absolutely for real part greater than $1$.
Let $\sigma_c\leq1$ be its abscissa of convergence; the case
$\sigma_c=-\infty$ is permitted.  For real part greater than $1$,
\begin{equation}\label{mellin:H-transform}
 F(s):=\int_1^\infty H(x)x^{-s-1}\,dx
   =\frac{P'(s)/P(s)-(a+bs)}{s(s-1)}+\frac C{s-\theta}.
\end{equation}
Denote the meromorphic expression on the right by $\Phi(s)$.
It is meromorphic for $\operatorname{Re}s>0$.
Equation~\eqref{mellin:P-at-one} shows that its singularity at $1$ is
removable; here $\theta<1$ and $P(1)\neq0$ are used.

We claim that $\sigma_c\leq\theta$.  Otherwise
$\theta<\sigma_c\leq1$ is a finite real number, and the identity theorem
for meromorphic functions gives $F=\Phi$ throughout
$\operatorname{Re}s>\sigma_c$.  If $\sigma_c=1$, the removability just
proved contradicts Lemma~\ref{mellin:landau}.
Suppose therefore that $\beta=\sigma_c\in(\theta,1)$.
If $P(\beta)\neq0$, then $\Phi$ is holomorphic near $\beta$, again
contradicting that lemma.
If $P(\beta)=0$, let $m\geq1$ be its order of vanishing.  Locally,
$P(s)=(s-\beta)^mU(s)$ with $U(\beta)\neq0$, so
\[
 \frac{P'(s)}{P(s)}=\frac m{s-\beta}+\frac{U'(s)}{U(s)}.
\]
Thus the residue of $\Phi$ at $\beta$ is
\[
 \frac{m}{\beta(\beta-1)}<0.
\]
As real $\sigma$ decreases to $\beta$ from above,
$F(\sigma)=\Phi(\sigma)\to-\infty$.  This is impossible because
$F(\sigma)$ is the integral of a nonnegative real function.
The claim follows.

Consequently $F$ is holomorphic in $\operatorname{Re}s>\theta$.
Define there the holomorphic coefficient
\[
 Q(s)=a+bs+s(s-1)\left(F(s)-\frac C{s-\theta}\right).
\]
Multiplication of~\eqref{mellin:H-transform} by $s(s-1)P(s)$ gives
$P'=QP$ for real part greater than $1$.  Both sides are holomorphic
in the connected half-plane $\operatorname{Re}s>\theta$, so the
identity theorem extends this equality throughout that half-plane.
A nonzero holomorphic function satisfying $P'=QP$ cannot have a zero:
at a zero of order $m\geq1$, the order of $P'$ would be $m-1$, whereas
the order of $QP$ would be at least $m$.  Hence $P$ has no zeros there.
Every exponent in its product is a positive integer.  It follows that
each $L(s,\chi)$ and $Z(s)$ is nonzero there.  Finally
$Z(s)=(s-1)\zeta(s)$ away from $1$, proving the assertion for $\zeta$.
\end{proof}

The argument uses no estimate for $-B(x)$.
Its essential sign is the negative residue
$m/(\beta(\beta-1))$ for $0<\beta<1$.  Positive multiplicities prevent
cancellation of zero orders in the product $P$.

\subsection{An absolute-bound criterion and the limitation of windows}

\begin{corollary}\label{mellin:absolute}
Fix $\theta\geq0$.  If, for every $\epsilon>0$,
$G(x)=O_\epsilon(x^{\theta+\epsilon})$ in absolute value, then
$\zeta$ has no zeros with real part greater than $\theta$.
For a fixed nonprincipal character with $L(1,\chi)\neq0$, if
$g_\chi(x)=O_\epsilon(x^{\theta+\epsilon})$ in absolute value for
every $\epsilon>0$, then $L(s,\chi)$ has no zeros in that half-plane.
\end{corollary}

\begin{proof}
The assumptions make the respective Mellin integrals $I(s)$ and $J(s)$
holomorphic in $\operatorname{Re}s>\theta$.  Indeed, on any compact
subset choose $\epsilon$ smaller than its distance from the boundary;
the resulting power bound dominates the integrands and all their
derivatives locally uniformly.  The identities of
Lemma~\ref{mellin:identities}, multiplied out, become
\[
 Z'(s)=\bigl[-1+(1+\gamma)s-s(s-1)I(s)\bigr]Z(s),
\]
and
\[
 L'(s,\chi)=\bigl[c_\chi s-s(s-1)J(s)\bigr]L(s,\chi).
\]
They initially hold for real part greater than $\max(1,\theta)$ and
extend to $\operatorname{Re}s>\theta$ by the identity theorem.
The same order-of-vanishing argument as above excludes zeros.
The elementary continuations are holomorphic throughout this
half-plane because $\theta\geq0$.
\end{proof}

The global quantifier in~\eqref{mellin:global-bound} must be preserved.
An estimate only on arbitrarily large intervals
$[Y,Y^\lambda]$, with fixed $\lambda>1$, does not by itself make
$Cx^\theta-B(x)$ nonnegative on the entire integration range.  It
therefore cannot be substituted into the proof of
Theorem~\ref{mellin:one-sided-family}.

For a general continuous function, even an absolute bound on such
intervals gives no global growth estimate.  To see this, choose $Y_1>1$ and numbers
$Y_j\to\infty$ so rapidly that
$Y_{j+1}>4Y_j^\lambda$.  Let $\eta$ be a continuous nonnegative function
supported on $[2,3]$, with $\eta(5/2)=1$, and fix $\beta>\theta$.
The locally finite sum
\[
 f(x)=\sum_{j\geq1}x^\beta\eta(x/Y_j^\lambda)
\]
vanishes on every $[Y_j,Y_j^\lambda]$, while
$f((5/2)Y_j^\lambda)=((5/2)Y_j^\lambda)^\beta$.
It is not $O(x^\theta)$.  Moreover, if $\eta$ is positive on a
subinterval of $(2,3)$, its Mellin integral diverges for real exponents
at most $\beta$: each separate bump contributes a positive constant
times $(Y_j^\lambda)^{\beta-\sigma}$.

This example is not a counterexample to a theorem about the particular
prime sums above.  Their logarithmic-derivative structure may allow a
stronger theorem from long multiplicative intervals.  Such a theorem
requires an additional argument.  The global criterion proved here
does not establish it, and no global estimate is inferred from a
window estimate in this section.
